\section{Conclusion}
\label{sec:conclusion}

We introduced Continuous Interaction Diffusion, a diffusion-native model--runtime architecture that replaces isolated tool-call rounds with persistent asynchronous perception over revisable cognition. CID separates externally controlled facts, typed continuous thought, and user-visible display; exposes information needs through persistent bindings; reuses or refreshes external observations according to source semantics; and allows arriving evidence to reopen affected state during denoising. Our evaluation is centered on the resulting end-to-end questions: whether the corrected trained model uses tools reliably, whether external latency can be overlapped with useful computation, and whether delayed or changing evidence can revise existing cognition without unnecessary restart or repeated I/O. Together, these tests make the central CID claim directly falsifiable against simpler explicit-call and asynchronous autoregressive alternatives.
